\documentclass[12pt]{article}
\usepackage[utf8]{inputenc}
\usepackage[T1]{fontenc}
\usepackage{hyperref}
\title{Western Minyag Language Resources for K-12 Curriculum Development}
\author{Compiled by Claude Sonnet 5.5 (claude-sonnet-5-5) for Andr\'as Kornai}
\date{}
\begin{document}
\maketitle

\section{Introduction}
Western Minyag (ISO 639-3 wmg; Glottocode west2417) is a Sino-Tibetan (Qiangic) language of western Sichuan, China. Sources disagree on speaker numbers and locations. One source gives about 10,000 speakers in villages around Gangkar (Gongga) Mountain in Kangding and Jiulong Counties. Another gives about 11,500, and Joshua Project gives 11,000. The language is not known to be taught in schools. Little pedagogical material appears to exist; the search results showed mainly documentation and reference material.

\section{Primary Educational Resources}
No primers, textbooks or school materials in Western Minyag were found in search results.

\section{Linguistic and Reference Materials}
\begin{itemize}
\item Wikipedia, Muya language (covers the Minyag varieties): \url{https://en.wikipedia.org/wiki/Muya_language}
\item ELDP/ELAR project profile, ``Grammar of Minyag, a minority language of Western Sichuan'': \url{https://blogs.soas.ac.uk/elar/2020/07/10/eldp-project-profile-grammar-of-minyag-a-minority-language-of-western-sichuan/}
\item Numerals resource with language name and location for Minyag (Western): \url{https://lingweb.eva.mpg.de/channumerals/Minyag-Western.htm}
\end{itemize}

\section{Religious and Community Resources}
Joshua Project reports that Bible translation has started. This is a people-group profile, not a scripture text: \url{https://joshuaproject.net/people_groups/18599/CH}. The Asia Harvest people-group profile is also listed, but it is only a profile: \url{https://people-groups.asiaharvest.org/people-group-profiles/sichuan/minyak.pdf}

\section{Digital and Audio Resources}
The ELDP archive, ``Western Minyag, an endangered language of Western Sichuan'', holds an audio-visual corpus. Its annotations are reported to cover over 10 hours of data as ELAN files, plus a FLEx database: \url{https://eldp.access.preservica.com/dk0539/}

\section{Curriculum Development Context}
Western Minyag is spoken by all adults but not all young people, and it is not known to be taught in schools. A curriculum would have to be built from scratch on the documentation corpus and grammar work above. That would need collaboration with the researchers and the community. Dialect labels and codes vary between sources (Western Minyag is wmg, Eastern Minyag is emq), so they must be checked before use. No orthography for school use was found.

\end{document}